\begin{table}[htbp]
\centering\small
\begin{tabular}{lllrrrrr}
\toprule
Comparison & Dataset & Task & $n$ & $\pm .005$ & $\pm .010$ & $\pm .020$ & Both signs$^*$ \\
\midrule
8--32 & APT & coarse & 2 & 0 & 0 & 0 & 1 \\
8--32 & APT & fine & 2 & 0 & 0 & 2 & 0 \\
8--32 & COMBAT & coarse & 2 & 0 & 2 & 2 & 0 \\
8--32 & COMBAT & fine & 2 & 0 & 1 & 2 & 0 \\
8--32 & OneK1K & coarse & 2 & 0 & 0 & 2 & 0 \\
8--32 & OneK1K & fine & 2 & 0 & 1 & 2 & 0 \\
Extended LR & COMBAT & coarse & 2 & 2 & 2 & 2 & 0 \\
Extended LR & COMBAT & fine & 2 & 0 & 2 & 2 & 0 \\
Extended LR & OneK1K & coarse & 4 & 2 & 3 & 4 & 0 \\
Extended LR & OneK1K & fine & 4 & 0 & 4 & 4 & 0 \\
\bottomrule
\end{tabular}
\caption{Post-hoc practical-equivalence sensitivity: number of released joint 95\% intervals strictly inside each stated margin. The 8--32 rows use $T=6{,}400$ and both LR and XGBoost; extended rows use LR at both totals for adjacent subject-budget comparisons. $^*$Intervals reaching both $-0.01$ and $+0.01$, indicating insufficient precision to exclude meaningful effects in either direction. Counts are pointwise diagnostics, not formal TOST decisions or evidence that an entire scaling frontier is equivalent.}
\label{tab:equivalence_sensitivity}
\end{table}
